\documentclass[11pt]{article}
\usepackage[utf8]{inputenc}
\usepackage{amsmath,amssymb,amsthm}
\usepackage[margin=1in]{geometry}
\usepackage{hyperref}
\usepackage{xcolor}
\usepackage{enumitem}
\usepackage{newunicodechar}
\newunicodechar{ℝ}{\ensuremath{\mathbb{R}}}
\newunicodechar{⟨}{\ensuremath{\langle}}
\newunicodechar{⟩}{\ensuremath{\rangle}}
\newunicodechar{α}{\ensuremath{\alpha}}
\newunicodechar{β}{\ensuremath{\beta}}
\newunicodechar{γ}{\ensuremath{\gamma}}
\newunicodechar{λ}{\ensuremath{\lambda}}
\newunicodechar{μ}{\ensuremath{\mu}}
\newunicodechar{σ}{\ensuremath{\sigma}}
\newunicodechar{φ}{\ensuremath{\varphi}}
\newunicodechar{ψ}{\ensuremath{\psi}}
\newunicodechar{κ}{\ensuremath{\kappa}}

\title{Tide retrospective:\\\emph{perturbed harmonic-Gibbs moments
$\langle x^k\rangle_h$ for $k = 3, 4$ (O1)}}
\author{Daniel Murfet}
\date{7 May 2026}

\begin{document}
\maketitle

\begin{abstract}
This retrospective records a 528-line programme continuation of I4
(global-constancy of the harmonic Gibbs covariance, merged earlier
today). I4 closed $k = 1, 2$: $\langle x\rangle_h = -h/\lambda$ and
$\langle x^2\rangle_h = 1/(\lambda t) + h^2/\lambda^2$. This tide
extends to $k = 3, 4$ via the same square-completion shift method,
and lands the **public Gibbs-expectation transport lemma** GPT-5.5~Pro
named at the deliberation: $\langle f(x)\rangle_h = \langle
f(y - h/\lambda)\rangle_0$ for any observable $f$. The decisive
mathematical observation: the normalised perturbed law is exactly
$\mathcal N(-h/\lambda,\,1/(\lambda t))$, so the closed forms are
the raw Gaussian moments. Single-session run, build green after three
small fixes (a cast normalisation for $((2\!\cdot\!2-1)!!)$ in
$\int y^4$, an rpow-algebra sequencing in $\int y^4$, and the standard
\texttt{Pi.add}-vs-single-lambda workaround per the laplace
\texttt{CLAUDE.md}). Zero \texttt{sorry}, zero \texttt{axiom},
zero \texttt{native\_decide}.
\end{abstract}

\section{Setting}

I4 (global-constancy harmonic, merged this morning) proved that the
harmonic Gibbs covariance is globally constant in $h$:
$\mathrm{Cov}_h(x, x) = 1/(\lambda t)$ for all $h \in \mathbb R$. Its
proof routed through closed-form evaluations of three perturbed
integrals via a square-completion shift identity. I4's retrospective
\S Follow-ups~5 explicitly flagged: \emph{``Promote
\texttt{integral\_with\_perturbation\_eq\_shifted} on second use.''}

This O1 tide is that second use. The user's direction was the family
$\langle x^k\rangle_h$ for $k \in \{1, 2, 3, 4\}$, with explicit
closed forms named by binomial expansion of $(y - h/\lambda)^k$. The
deliberation converged on Candidate A+ (specific theorems for
$k = 3, 4$ \emph{plus} a public shift-at-expectation transport lemma)
rather than B (a parametric general-$k$ statement) or C (just $k = 3$).
Both Claude and GPT-5.5~Pro voted A+ on the first round.

\section{The thing formalised}

\begin{quote}\itshape
For the harmonic Gibbs measure
$\mu_{t,h}(x) \propto \exp(-t((\lambda/2)x^2 + h\cdot x))$ on
$\mathbb R$ against Lebesgue with $\lambda, t > 0$, the perturbed
moments are
\[
  \langle x^3\rangle_h = -\frac{3h}{\lambda^2 t} - \frac{h^3}{\lambda^3},
  \qquad
  \langle x^4\rangle_h = \frac{3}{\lambda^2 t^2}
    + \frac{6h^2}{\lambda^3 t} + \frac{h^4}{\lambda^4}.
\]
\end{quote}

\noindent The Lean signatures (public):
\begin{verbatim}
theorem gibbsExp_harmonic_h_eq_unperturbed_shift
    {lam t : ℝ} (hlam : 0 < lam) (ht : 0 < t)
    (h : ℝ) (f : ℝ → ℝ) :
    Threepoint.gibbsExp (volume : Measure ℝ)
        (fun x ↦ lam / 2 * x ^ 2)
        (fun x ↦ x) t h f
    = Threepoint.gibbsExp (volume : Measure ℝ)
        (fun x ↦ lam / 2 * x ^ 2)
        (fun x ↦ x) t 0
        (fun y ↦ f (y - h / lam))

theorem gibbsExp_h_cube_harmonic_eq
    {lam t : ℝ} (hlam : 0 < lam) (ht : 0 < t) (h : ℝ) :
    Threepoint.gibbsExp ... t h (fun x ↦ x ^ 3)
    = -3 * h / (lam ^ 2 * t) - h ^ 3 / lam ^ 3

theorem gibbsExp_h_quartic_harmonic_eq
    {lam t : ℝ} (hlam : 0 < lam) (ht : 0 < t) (h : ℝ) :
    Threepoint.gibbsExp ... t h (fun x ↦ x ^ 4)
    = 3 / (lam ^ 2 * t ^ 2) + 6 * h ^ 2 / (lam ^ 3 * t)
        + h ^ 4 / lam ^ 4
\end{verbatim}

\section{Proof strategy}

The decisive observation, due to GPT-5.5~Pro: the normalised
perturbed law is exactly Gaussian with mean $-h/\lambda$ and variance
$1/(\lambda t)$. So the moments are the standard raw Gaussian moments
$E[X^k] = \mu^k + \binom{k}{2}\mu^{k-2}\sigma^2 + \binom{k}{4}
\mu^{k-4}\sigma^4 + \cdots$ (only even-power-of-$\sigma$ terms
survive by parity). Substitute $\mu = -h/\lambda$, $\sigma^2 =
1/(\lambda t)$ and the closed forms drop out.

The Lean proof assembles in three layers.

\paragraph{Layer 1: re-derive the I4 shift identity locally.} The
integral-level shift identity is \texttt{private} in the I4 file. To
keep this tide self-contained, I re-prove it under a primed name. The
proof is verbatim from I4: square-completion gives the pointwise
identity $\exp(-t(L_h(x))) = \exp(t h^2/(2\lambda)) \cdot
\exp(-t(L_0(x + h/\lambda)))$, which integrates to
$\int f(x) e^{-tL_h} dx = e^{th^2/(2\lambda)}\int f(y - h/\lambda)
e^{-tL_0} dy$ via \texttt{integral\_add\_right\_eq\_self} for the
shift. About 30 lines.

\paragraph{Layer 2: integrability witnesses + closed-form moments.}
Build a single $\int y^n e^{-t(\lambda/2)y^2}$ integrability lemma
for any $n$ via Mathlib's polynomial-times-Gaussian integrability
result, after a coefficient-matching rewrite.\footnote{Mathlib lemma:
\texttt{integrable\_pow\_mul\_exp\_neg\_mul\_sq}.} Build closed forms
for $n = 0, 1, 2, 3, 4$ by specialising Tide~10's even/odd
moments.\footnote{\texttt{harmonic\_int\_pow\_even} and
\texttt{harmonic\_int\_pow\_odd}.} $n$ odd gives 0; $n = 0$ gives
$\sqrt{2\pi/(\lambda t)}$; $n = 2$ gives $(1/(\lambda t))
\sqrt{2\pi/(\lambda t)}$ via $(\lambda t)^{-(1+1/2)} = (\lambda t)^{-1}
\cdot (\lambda t)^{-1/2}$; $n = 4$ gives $(3/(\lambda t)^2)
\sqrt{2\pi/(\lambda t)}$ via the same rpow-additivity argument with
$(\lambda t)^{-(2+1/2)}$. About 100 lines for the five bridges.

\paragraph{Layer 3: the transport lemma and the two moments.} The
transport lemma's proof is short: apply the shift identity to both
numerator and denominator of $\langle f\rangle_h$. Each carries an
$\exp(t h^2/(2\lambda))$ factor; \texttt{field\_simp} cancels them in
the ratio, delivering the unperturbed expectation of the shifted
observable. The cube and quartic theorems then expand $(y - h/\lambda)^k$
pointwise via \texttt{funext y; ring}, split the integral via
\texttt{integral\_add} and \texttt{integral\_const\_mul}, plug in the
$n = 0, 1, 2, 3, 4$ closed forms, and finish with
\texttt{field\_simp} (no trailing \texttt{ring} needed —
\texttt{field\_simp} closes the rational-function identity outright).
About 90 lines per moment.

\section{Roadblocks and resolutions}

Three small fixes between draft and green build, none load-bearing.

\paragraph{Cast normalisation in $\int y^4$.} The
\texttt{harmonic\_int\_pow\_even} statement carries
$((2k-1)\!\!)\!\!\,!! : \mathbb R$ as a coefficient. At $k = 2$ this
is $((2\cdot 2 - 1)\!\!)\!\!\,!! = 3!! = 3$, but the cast needed
explicit unfolding (a \texttt{change} to $((3 : \mathbb N)!! :
\mathbb R)$, then \texttt{decide} on the natural-number doubleFactorial,
then \texttt{norm\_num} on the cast). The $((2 : \mathbb N) : \mathbb R)
= 2$ exponent normalisation likewise needed an explicit
\texttt{norm\_cast}. About 5 lines of cast plumbing per even-power
closed form.

\paragraph{rpow algebra in $\int y^4$.} The closed form $(\lambda t)^{-(2 +
1/2)} = (\lambda t)^{-2} \cdot (\lambda t)^{-1/2}$ — straightforward
mathematically — is fragile in Lean because the parenthesisation of
\texttt{rpow\_add} matters and the goal can have several subterms
matching the rewrite pattern. First attempt with a chained rewrite
sequence failed at \emph{``Tactic rewrite failed: did not find an
occurrence of the pattern''} on the third step. Fix: factor out a
single named \texttt{hcombine} hypothesis with \texttt{Real.rpow\_add}
inside, then $\texttt{rw [hcombine]; ring}$ closes the goal in one
shot. The lesson: when chaining rpow rewrites, the named-intermediate
form is more robust than the inline-show pattern.

\paragraph{\texttt{Pi.add} vs single-lambda mismatch in
\texttt{integral\_add}.} Standard issue per the laplace
\texttt{CLAUDE.md}: \texttt{Integrable.add} returns
\texttt{Integrable (f + g)} where the sum is \texttt{Pi.add}, but
\texttt{integral\_add}'s rewrite pattern is the single-lambda form
\texttt{fun~y =>~f y + g y}. Lean's pattern-matcher doesn't unfold
\texttt{Pi.add} automatically, so the rewrite fails on the elaborated
goal. Fix: introduce a type-ascribed single-lambda integrability
witness \texttt{have h\_combined : Integrable (fun y => ...) volume :=
h\_left.add h\_right} before the \texttt{rw [integral\_add ...]}. This
issue surfaced in both the cube and quartic numerator integrals; the
fix added six \texttt{have} blocks across the two theorems.

\paragraph{Trailing \texttt{ring} after \texttt{field\_simp}.} Both
the cube and quartic theorems' final lines originally ended with
\texttt{field\_simp; ring}. After fixing the \texttt{Pi.add} issue, the
\texttt{field\_simp} step closed the goal completely, leaving no
goals for the trailing \texttt{ring}. Removing the dead \texttt{ring}
fixed two \emph{``No goals to be solved''} errors.

\section{What was Mathlib, what was new}

\paragraph{Mathlib.} The Bochner-integral linearity lemmas
(translation invariance, additivity, scalar multiplication, congruence
under a.e.\ equality), the polynomial-times-Gaussian integrability
result, and the standard \texttt{rpow} algebra (additivity in the
exponent, negation, sqrt-as-rpow, Nat-cast).\footnote{Names:
\texttt{integral\_add\_right\_eq\_self}, \texttt{integral\_const\_mul},
\texttt{integral\_add}, \texttt{integral\_sub},
\texttt{integral\_congr\_ae},
\texttt{integrable\_pow\_mul\_exp\_neg\_mul\_sq},
\texttt{Real.rpow\_add}, \texttt{Real.rpow\_neg},
\texttt{Real.sqrt\_eq\_rpow}, \texttt{Real.rpow\_natCast}.}
Plus \texttt{field\_simp} and \texttt{ring} as closing tactics. All
familiar; no new Mathlib API surfaced.

\paragraph{Laplace seabed (existing).} I4's integral-level shift
identity (re-derived locally, since the I4 original is
\texttt{private}). Tide~10's even/odd unperturbed-moment closed forms
for $n = 0, 1, 2, 3, 4$ and \texttt{partitionFunction\_harmonic}. The
standard \texttt{Threepoint.gibbsExp} definition.

\paragraph{New.} Five $\int y^n e^{-t(\lambda/2)y^2}$ closed-form
bridges (named \texttt{int\_pow\_zero/one/two/three/four}), one
general polynomial-times-Gaussian integrability witness
parameterised by $n : \mathbb N$, the public shift-at-expectation
transport lemma, and the two public moment theorems. The transport
lemma is the load-bearing new public API: it captures the
perturbed-law-equals-shifted-unperturbed-law fact in a form
independent of the choice of observable, and is the natural lift of
I4's integral-level shift identity to the expectation level.

\section{Lessons}

\begin{itemize}
\item \emph{The transport lemma was the right abstraction.} GPT-5.5~Pro
  named it at the deliberation, replacing my draft's plan to promote
  the integral-level shift identity. The transport lemma's stronger
  invariant — \emph{the perturbed law is the shifted unperturbed law,
  full stop} — yields shorter cube/quartic proofs (no need to manage
  the $\exp(t h^2/(2\lambda))$ factor explicitly at every site;
  \texttt{field\_simp} cancels it once inside the transport lemma's
  proof).
\item \emph{Re-deriving I4's private helpers is cheaper than
  un-privating.} The integral-level shift identity is $\sim$30 lines.
  Modifying the I4 file to expose it would have crossed the
  ``edit-merged-file'' boundary and risked merge surprises with the
  L2 in-flight tide that branched off N1 (which itself was just
  merged). Local re-derivation kept the change contained.
\item \emph{The \texttt{Pi.add} workaround is a recurring pattern.}
  This is the third tide today (after I4 itself and G2+) where the
  \texttt{Integrable.add}-returns-\texttt{Pi.add} vs.\
  \texttt{integral\_add}-expects-single-lambda mismatch surfaced. The
  laplace \texttt{CLAUDE.md} already documents the workaround; future
  tides should reach for the type-ascribed-single-lambda witness on
  first sight, not after a failed rewrite.
\item \emph{One parametric integrability witness beats three
  specialised ones.} I4 used three separate witnesses for the
  unperturbed Gaussian moments (zeroth, first, second power). This
  tide uses one parametric witness taking $n : \mathbb N$, derived
  from Mathlib's polynomial-times-Gaussian integrability via a
  coefficient rewrite. Worth promoting to lean-common when a third
  consumer materialises.
\end{itemize}

\section{Follow-ups}

\begin{itemize}
\item \emph{Generic-$k$ moments via the binomial sum.} Given the
  parametric integrability lemma above, a single
  \texttt{gibbsExp\_h\_pow\_harmonic\_eq} for any $k : \mathbb N$
  (Candidate B from this tide's deliberation, deferred) would land in
  $\sim$80 additional lines: $\sum_{j=0}^{k} \binom{k}{j}
  (-h/\lambda)^{k-j} \langle y^j\rangle_0$, with a separate
  \texttt{gibbsExpectation\_harmonic\_pow\_*\_y} lemma for general
  $\langle y^j\rangle_0$. Defer until a $k = 5$ or $k = 6$ consumer
  arrives.
\item \emph{Translation invariance for arbitrary potentials.} The
  shift identity here is harmonic-specific (square completion). A
  generic potential $L$ admits a transport lemma of the form
  $\langle f(x)\rangle_{L, h} = \langle f(y - \tau(h))\rangle_{L, 0}$
  only when $L$ has appropriate symmetry. Worth recording as a
  \emph{not-tide-shaped} candidate: depending on the symmetry hypothesis
  bundle, this could be a 200-line expansion or a multi-tide
  programme.
\item \emph{\texttt{\textbackslash leanref}-tag the primer.} The
  primer's discussion of the shifted-Gaussian centroid mentions
  $\langle x^3\rangle_h$ and $\langle x^4\rangle_h$ as exercises in
  perturbative expansion. Pin
  \texttt{\textbackslash laplaceCommit = ea3e13e} (this tide's commit)
  and tag both. Pairs with K2 / N4 in the survey.
\item \emph{Discharge mixed-power \texttt{GibbsObservable} hypotheses
  (O2 in the survey).} The transport lemma and the cube/quartic moments
  may discharge some of the mixed-power \texttt{GibbsObservable}
  hypotheses Tide~13 left conditional. Worth a follow-up audit.
\end{itemize}

\section*{Acknowledgements}

GPT-5.5~Pro's deliberation was decisive on three counts. First, the
recasting of the closed-form derivation in terms of Gaussian raw
moments $E[X^k] = \mu^k + \binom{k}{2}\mu^{k-2}\sigma^2 + \cdots$
under $X \sim \mathcal N(-h/\lambda,\,1/(\lambda t))$ replaced the
binomial-of-$(y - h/\lambda)^k$ derivation in my draft with the
shorter Gaussian-moment derivation. Second, the recommendation to
land a public Gibbs-expectation transport lemma rather than just the
integral-level shift identity. Third, a numerical-sanity table of
five $(\lambda, h, t)$ cases and their closed-form values for
verification (which matched scipy.quad to machine precision). Full
consult preserved at the GPT v1 file.\footnote{Preserved at
\texttt{projects/primer/tide-log/gpt55\_tide\_perturbed\_moments\_harmonic\_v1.md}.}

This Tide branched off the post-N1-merge \texttt{main} (commit
\texttt{47e3426}). I4 is at \texttt{5aa212c}; the transport lemma
re-derives I4's private shift identity locally rather than modifying
the I4 file. No in-flight tides interact with this work (L2 / Q1 / A2
all on parallel sessions, all independent of the perturbed harmonic
moments).

\end{document}
